%%%%%%%%%%%%%%%%%%%%%%%%%%%%%%%%%%%%%%%%%%%%%%%%%
%%%%%%%%%%%%%%%%%%% EXAMPLE 2 %%%%%%%%%%%%%%%%%%%
%%%%%%%%%%%%%%%%%%%%%%%%%%%%%%%%%%%%%%%%%%%%%%%%%

% Minimal example of an HDR demand, i.e., the short document sent to the university
% to be allowed to register for the HDR -- here written in French
% Compile with: latexmk -lualatex -outdir=build main_demand.tex
% Compared with "main_hdr.tex", the "hdr_demand" option gives a shorter document:
% no table of contents, no chapter pages, no back cover, tighter spacing

%%%%%%%%%%%%%%%%%%%%%%%%%%%%%%%%%%%%%%%%%%%%%%%%%
%%%%%%%%%%%%%%%%% CONFIGURATION %%%%%%%%%%%%%%%%%
%%%%%%%%%%%%%%%%%%%%%%%%%%%%%%%%%%%%%%%%%%%%%%%%%

% Choose template
\documentclass[hdr_demand, french]{template/hdr}

%%%%%%%%%%%%%%%%%%%%%%%%%%%%%%%%%%%%%%%%%%%%%%%%%
%%%%%%%%%%%%%%%%%% INFORMATION %%%%%%%%%%%%%%%%%%
%%%%%%%%%%%%%%%%%%%%%%%%%%%%%%%%%%%%%%%%%%%%%%%%%

% Author
\setAuthorFirstName{Prénom}
\setAuthorName{Nom}

% Affiliations
\setUniversity{Université Fictive}
\setDoctoralSchool{\acroname{ed} Fictive}
\setDoctoralSchoolSubtitle{Sciences et Technologies de l'Information}
\addLogo{manuscript/covers/logo_university.png}{https://example.org}
\addLogo{manuscript/covers/logo_doctoral_school.png}{https://example.org}
\addLogo{manuscript/covers/logo_laboratory.png}{https://example.org}

% Document
\setDocumentTitle{Un titre fictif pour un document fictif}

% Testimony
% Specific to the HDR demand: a sworn statement, followed by your signature
\setSignature{manuscript/covers/signature.png}
\setTestimony{
    Je soussigné·e, \authorFirstName~\textsc{\authorName}, déclare sur l'honneur n'avoir pas déposé de demande d'inscription en \acrofull{hdr} dans un autre établissement que l'Université Fictive. \\
    Ce paragraphe est un texte de remplacement : remplacez-le par les motivations réelles de votre demande, par exemple le rattachement de votre laboratoire, la proximité géographique, ou vos collaborations en cours.}

% Bibliography
% Fewer categories than in the full manuscript, to keep the document short
% Every category cited in the text must be declared here, otherwise compilation fails
% Categories that are declared nowhere simply do not appear (here, "book")
\addBiblioCategory{thesis}{T}{Thèses}
\addBiblioCategory{journal}{J}{Journaux}
\addBiblioCategory{conference}{C}{Conférences}
\addBiblioCategory{workshop}{W}{Workshops}
\addBiblioCategory{preprint}{P}{Preprints}

%%%%%%%%%%%%%%%%%%%%%%%%%%%%%%%%%%%%%%%%%%%%%%%%%
%%%%%%%%%%%%%%%%%%%% CONTENT %%%%%%%%%%%%%%%%%%%%
%%%%%%%%%%%%%%%%%%%%%%%%%%%%%%%%%%%%%%%%%%%%%%%%%

% Go
\begin{document}

% Chapters
% In a short document, there is no dedicated bibliography chapter:
% \printbibliography is called at the end of "research.tex", as a section of that chapter
% Only the categories declared above are printed, external references and online
% resources are cited inline but not listed
% Unnumbered chapter, obtained with the starred variant of \chapter
% It works in the short document too, where chapters are much lighter
\chapter*[demand_foreword]{Avant-propos}

Ce chapitre n'est pas numéroté : il est obtenu avec \verb|\chapter*[label]{Titre}|.
Comme les autres, il porte une étiquette, et peut donc être cité avec \verb|\nameref| : \verb|\nameref{demand_foreword}| donne \nameref{demand_foreword}.

Lorem ipsum dolor sit amet, consectetur adipiscing elit, sed do eiusmod tempor incididunt ut labore et dolore magna aliqua.
Ut enim ad minim veniam, quis nostrud exercitation ullamco laboris nisi ut aliquip ex ea commodo consequat.

% In a short document (option "hdr_demand"), chapters are much lighter:
% no dedicated chapter page, no background image, no mini table of contents
\chapter[demand_cv]{Curriculum Vitæ}
 [
  Ce chapitre présente les informations générales de mon dossier, et résume mon parcours, mes encadrements et mes enseignements. \\
  Afin d'alléger le document, la plupart des acronymes ne sont pas explicités. La version \acro{pdf} de ce document permet toutefois d'afficher leur nom complet au survol de la souris\footnote{Nécessite un lecteur \acro{pdf} compatible avec les \emph{tooltips} (\eg Acrobat, Evince, Overleaf).}. Essayez avec ce paragraphe !
 ]

\section[demand_cv_info]{Informations générales}

\begin{tabularx}{\linewidth}{r|X}
    \textbf{Nom}         & Prénom \textsc{Nom}                                                   \\
    \textbf{Né·e le}     & \date{1988-10-03}                                                     \\
    \textbf{Nationalité} & Fictive                                                               \\
    \textbf{Liens}       & \href{https://example.org}{\faGlobeAmericas}
    \hspace{\valparskip}\href{https://example.org}{\faGraduationCap}
    \hspace{\valparskip}\href{https://example.org}{\faGithub}                                    \\
    \textbf{Email}       & \href{mailto:prenom.nom@example.org}{\texttt{prenom.nom@example.org}}
\end{tabularx}

\section[demand_cv_parcours]{Parcours académique}

\begin{tabularx}{\linewidth}{r|X}
    \textbf{2010 -- 2015} & Doctorat en sciences fictives \newline
    Laboratoire Fictif, Université Fictive \newline
    Thèse : \cite{lastname2015thesis}                              \\
    \textbf{2015 -- 2017} & Post-doctorat \newline
    Institut Fictif                                                \\
    \textbf{Depuis 2017}  & Maître de conférences \newline
    Laboratoire Fictif, Université Fictive
\end{tabularx}

\section[demand_cv_encadrement]{Encadrements}

\begin{tabularx}{\linewidth}{r|X}
    \textbf{2021 -- 2024} & Fourth \textsc{Student}, doctorant·e (50\%) \newline
    \inquotes{Un titre fictif pour une thèse fictive} \newline
    Publications : \cite{lastname2023workshop}, \cite{lastname2025preprint}                \\
    \textbf{2022}         & Second \textsc{Coauthor}, stagiaire de master (100\%) \newline
    \inquotes{Un autre sujet fictif}
\end{tabularx}

\section[demand_cv_enseignement]{Enseignements}

\begin{tabularx}{\linewidth}{r|X}
    \textbf{Depuis 2017} & \inquotes{Introduction à l'\acrofull{ai_fr}} \newline
    Niveau master 1, 30h/an                                                      \\
    \textbf{Depuis 2019} & \inquotes{Programmation fictive} \newline
    Niveau master 2, 45h/an
\end{tabularx}

\chapter[demand_research]{Activités de recherche}
 [
  Ce chapitre résume mes activités de recherche passées, présentes et futures. \\
  Tout le texte de ce document est un texte de remplacement : remplacez-le par le vôtre. \\
  Pour des raisons de place, les travaux tiers ne sont pas listés en fin de document. Vous trouverez donc çà et là des références à des articles~\cite{ipsum1998foundational} ou à des ressources en ligne~\cite{placeholder_www}, munies d'une icône cliquable menant à la page concernée. Comme pour les acronymes, vous pouvez survoler la référence pour en avoir un aperçu.
 ]

\section[demand_research_global]{Vision globale}

Lorem ipsum dolor sit amet, consectetur adipiscing elit, sed do eiusmod tempor incididunt ut labore et dolore magna aliqua.
Mes travaux portent sur l'\acrofull{ai_fr}, et plus particulièrement sur les objets fictifs~\cite{ipsum1998foundational}.

Ut enim ad minim veniam, quis nostrud exercitation ullamco laboris nisi ut aliquip ex ea commodo consequat~\cite{lastname2021journal}.
Duis aute irure dolor in reprehenderit in voluptate velit esse cillum dolore eu fugiat nulla pariatur~\cite{amet2017seminal}.

\begin{boxenv}[demand_research_question]{Question}[Formuler le problème en une phrase]
    Les boîtes fonctionnent de la même manière que dans le manuscrit complet.
    Le type est libre : ici \inquotes{Question}, ailleurs \inquotes{Défi} ou \inquotes{Définition}.
\end{boxenv}

\section[demand_research_results]{Principaux résultats}

Les résultats de \nameref{demand_research_figure} illustrent le comportement attendu d'un modèle fictif.
Le détail des travaux correspondants est disponible dans \cite{lastname2022conference} et \cite{lastname2024journal}.

\begin{figure}[demand_research_figure]{Un exemple de figure dans le document court. La numérotation des figures est continue sur tout le document, ce qui rend les renvois comme \texttt{\textbackslash nameref\{demand\_research\_figure\}} non ambigus.}
    % Example of a vector figure, meant to be included with \input{} inside a figure environment
% Such fragments are typically produced by a script, so that data and rendering stay in sync
% The tikzpicture environment is patched by the template to never overflow \linewidth

\begin{tikzpicture}

    \begin{axis}
        [
            height      = 5cm,
            width       = \linewidth,
            xlabel      = {Number of training epochs},
            ylabel      = {Placeholder metric},
            xmin        = 0, xmax = 20,
            ymin        = 0, ymax = 1,
            grid        = major,
            grid style  = {separatorColor},
            legend pos  = south east,
            legend cell align = left,
        ]

        \addplot[titlesColor, line width=1.5pt, mark=*, mark size=1.5pt] coordinates
            {
                (0, 0.10) (2, 0.38) (4, 0.55) (6, 0.66) (8, 0.73)
                (10, 0.78) (12, 0.82) (14, 0.84) (16, 0.86) (18, 0.87) (20, 0.87)
            };
        \addlegendentry{Placeholder model}

        \addplot[linksColor, line width=1.5pt, dashed, mark=square*, mark size=1.5pt] coordinates
            {
                (0, 0.10) (2, 0.30) (4, 0.44) (6, 0.53) (8, 0.59)
                (10, 0.63) (12, 0.66) (14, 0.68) (16, 0.69) (18, 0.70) (20, 0.70)
            };
        \addlegendentry{Placeholder baseline}

    \end{axis}

\end{tikzpicture}

\end{figure}

\begin{equation}[demand_research_equation]
    \mathcal{L}(\theta) = \frac{1}{N} \sum_{i = 1}^{N} \left\| f_\theta(\bm{x}_i) - \bm{y}_i \right\|_2^2 + \lambda \left\| \theta \right\|_2^2
\end{equation}

L'\nameref{demand_research_equation} donne la fonction de coût minimisée, où $\bm{x}$ est une entrée, $\bm{y}$ sa cible, et $\theta$ les paramètres du modèle.

\section[demand_research_future]{Perspectives}

Excepteur sint occaecat cupidatat non proident, sunt in culpa qui officia deserunt mollit anim id est laborum.
Ces perspectives prolongent \cite{lastname2025preprint}, et s'appuient sur les ressources décrites dans \cite{placeholder_www}.

\begin{itemize}
    \item Une première direction fictive, sur les données fictives.
    \item Une deuxième direction fictive, sur les modèles fictifs.
    \item Une troisième direction fictive, sur les applications fictives.
\end{itemize}

%%%%%%%%%%%%%%%%%%%%%%%%%%%%%%%%%%%%%%%%%%%%%%%%%
%%%%%%%%%%%%%%%%%% BIBLIOGRAPHY %%%%%%%%%%%%%%%%%
%%%%%%%%%%%%%%%%%%%%%%%%%%%%%%%%%%%%%%%%%%%%%%%%%

% In a short document, the list of publications is a section of the current chapter
\printbibliography


% Done
\end{document}
